\subsection{temporal\_fence\_t}\label{sub:temporal} % (fold)
Following the difficulties identified in \cite{ge2019}, it was clear that
hardware improvements are needed before time protection can be fully achieved.
The paper \cite{wistoff2023} introduce one such hardware improvement. First, it
demonstrates the presence of serious timing channels within an in-order RISC-V
core and confirms that the current ISA lacks mechanisms to close them reliably.

Following this, they present two versions of the \texttt{fence.t} instruction:
a basic flush that exclusively flushes the L1 data and instruction caches, the
TLBs, and the branch predictors. Through their evaluation they find that
further stateful components - such as the LFSR generating a pseudo-random index
sequence for the cache replacement policy and the round-robin memory arbiters
of the core - still leak data through the microarchitectural state. As such,
they also implement a Full Flush, which clears these components along with
those covered by the basic flush. The \texttt{fence.t} instruction proceeds as
follows:
\begin{itemize}
  \item Save the program counter.
  \item Save locally modified (dirty) architectural state.
  \item Drain pending transactions.
  \item Clear components that are not cleared on reset.
  \item Assert Microreset. Clearing all flip-flops containing non-architectural
    state.
\end{itemize}

A full explanation of each step is provided in \cite{wistoff2023}.

Given this implementation, they augmented the seL4 kernel to call the
\texttt{fence.t} instruction on every context switch, finding an overhead of
around 0.7\%. They note that this should be seen as a pessimistic result due to
the use of idle tasks, and that in a setting with high memory contention the
latency may be almost entirely hidden.

\subsubsection{With software support}
While \texttt{fence.t} was effective and inexpensive on simple in-order cores,
a follow-up paper \cite{wistoff2025} argues that out-of-order execution CPUs
introduce challenges for the implementation. They address these challenges by
presenting a software-supported temporal fence \texttt{fence.t.s} methodology.
They show that this software-supported version is able to close all on-core
timing channels without measurable hardware overhead, and with a minimal
average performance impact of 1.0\%.


% section temporal\_fence\_t (end)
